\section{Anchor Dossier Model}

\subsection{Purpose}

This generated note turns the strongest anchors into case files. Each dossier gathers occurrences, contexts, minimal alternations, component reuse, and lexical search gates. It is designed to move from structural and semantic compatibility toward controlled partial decipherment.

\subsection{Summary}

\begin{table}[htbp]
\centering
\begin{tabular}{p{0.58\textwidth}r}
\toprule
\textbf{Metric} & \textbf{Value} \\
\midrule
Anchor dossiers & 11 \\
Strong onomastic dossiers & 7 \\
Cross-frame dossiers & 4 \\
Anchor occurrence rows & 59 \\
Direct minimal tests & 49 \\
Component minimal tests & 277 \\
Anchor components & 15 \\
Root-like components & 2 \\
Prefix-like components & 3 \\
Suffix-like components & 4 \\
\bottomrule
\end{tabular}
\caption{Anchor dossier summary.}
\end{table}

\subsection{Top Dossiers}

\begin{table}[htbp]
\centering
\footnotesize
\begin{tabular}{lp{0.36\textwidth}p{0.28\textwidth}r}
\toprule
\textbf{Unit} & \textbf{Class} & \textbf{Context} & \textbf{Conf.} \\
\midrule
240-482 & institutional/admin label candidate & Site=Harappa & 0.966 \\
176 & emblem-linked name/title candidate & Site=Mohenjo-daro & 0.796 \\
590-390 & emblem-linked name/title candidate & Material=Steatite & 0.716 \\
235-222 & emblem-linked name/title candidate & Symbol=Bull & 0.705 \\
032-220 & emblem-linked name/title candidate & Symbol=Bull & 0.681 \\
240-176 & controlled lexical candidate & Material=Steatite & 0.666 \\
235-840-032 & controlled lexical candidate & SemanticFrame=PrimeEntitySealFormula & 0.665 \\
032 & suffix/title/object-ending candidate & Site=Chanhu-daro & 0.946 \\
240-740-090 & administrative object/formula candidate & Symbol=Cros & 0.935 \\
140 & administrative object/formula candidate & Symbol=Goat:4 & 0.778 \\
\bottomrule
\end{tabular}
\caption{Highest-priority anchor dossiers.}
\end{table}

\subsection{Component Probe}

\begin{table}[htbp]
\centering
\footnotesize
\begin{tabular}{llp{0.28\textwidth}p{0.34\textwidth}}
\toprule
\textbf{Sign} & \textbf{Variable} & \textbf{Units} & \textbf{Interpretation} \\
\midrule
032 & A07 & 032; 032-220; 235-840-032 & root-like component; appears independently and inside compounds \\
240 & A02 & 240-176; 240-482; 240-740-090 & prefix/title/classifier-like component \\
176 & R02 & 176; 240-176 & root-like component; appears independently and inside compounds \\
235 & M01 & 235-222; 235-840-032 & prefix/title/classifier-like component \\
060 & A06 & 820-060 & suffix/final qualifier candidate \\
090 & A10 & 240-740-090 & suffix/final qualifier candidate \\
140 & M07 & 140 & unresolved component \\
220 & A09 & 032-220 & suffix/final qualifier candidate \\
\bottomrule
\end{tabular}
\caption{Recurring components inside anchor units.}
\end{table}

\subsection{Working Reading Classes}

\begin{table}[htbp]
\centering
\footnotesize
\begin{tabular}{lp{0.72\textwidth}}
\toprule
\textbf{Unit} & \textbf{Working hypothesis} \\
\midrule
240-482 & 240 may mark institutional/admin class; 482 may be local entity stem \\
176 & root-like entity name/title component; compare with 240-176 \\
590-390 & compound name/title/place candidate with final formula-marker-like 390 \\
235-222 & 235 may be formula/title prefix; 222 may be entity stem \\
032-220 & 032 may qualify 220 or mark a subtype in bull-seal context \\
240-176 & 240 may qualify or classify root-like 176 \\
235-840-032 & multi-part entity/title formula sharing 235 and 032 with other anchors \\
032 & suffix/title/object-ending candidate \\
\bottomrule
\end{tabular}
\caption{Reading classes allowed by current evidence. These are not phonetic readings.}
\end{table}
